\documentclass[12pt]{article}
\usepackage[utf8]{inputenc}
\usepackage[T1]{fontenc}
\usepackage{hyperref}
\title{Adonara Language Resources for K-12 Curriculum Development}
\author{Compiled by Claude Sonnet 5.5 (claude-sonnet-5-5) for Andr\'as Kornai}
\date{}
\begin{document}
\maketitle

\section{Introduction}
Adonara (ISO 639-3 adr) is spoken on Adonara island in eastern Indonesia. It is usually treated as part of the wider Lamaholot cluster, so most scholarly material covers Lamaholot rather than Adonara alone.

\section{Primary Educational Resources}
No primers, textbooks or school programmes for Adonara turned up in search results.

\section{Linguistic and Reference Materials}
\begin{itemize}
\item Wikipedia, ``Adonara language'': \url{https://en.wikipedia.org/wiki/Adonara_language}
\item Wikipedia, ``Lamaholot language'': \url{https://en.wikipedia.org/wiki/Lamaholot_language}
\item Review essay, ``Trilingual dictionary of Lamaholot'' (on Pampus's Lamaholot--Indonesian--German dictionary; the repository record is access-restricted): \url{https://openresearch-repository.anu.edu.au/items/45033458-91ea-463a-8a50-02577ed24b84}
\item Boli, Ramly and Usman, ``Variasi Dialek Bahasa Lamaholot di Pulau Adonara, Kabupaten Flores Timur'', a study of dialect variation (in Indonesian): \url{https://ojs.unm.ac.id/ijses/article/view/22053}
\item ASJP Adonara wordlist: \url{https://asjp.clld.org/languages/ADONARA}
\item Austronesian Comparative Dictionary, Lamaholot entry: \url{https://acd.clld.org/languages/720}
\end{itemize}

\section{Religious and Community Resources}
None found.

\section{Digital and Audio Resources}
None found beyond the wordlist above.

\section{Curriculum Development Context}
There are about 100,000 speakers and no known school use. Reference descriptions exist only at the Lamaholot level, and no free online Adonara dictionary was found. Curriculum work would need new materials, built with local speakers and informed by the Lamaholot literature.
\end{document}